% !PRINTED 349
\weberpage{349}

%% [running head, not body text -- commented out 2026-08-07] Page furniture (running
%% head, folio, signature) is excluded from this edition by standing policy.
%% \noindent\S\ 95.\hfill Äquivalente Formen und Ideale im quadratischen Körper.\hfill 349

Es seien nun umgekehrt $\varphi = (a, b, c)$, $\varphi' = (a', b', c')$ zwei
\gesperrt{äquivalente Formen mit positiven ersten Koeffizienten} $a, a'$, und
\weq{6}{\omega = -\frac{b+\sqrt{\Delta}}{2a}, \quad \omega' = -\frac{b'+\sqrt{\Delta}}{2a'}\,.}
Dann gibt es eine lineare Substitution $\begin{pmatrix} \alpha, & \beta \\ \gamma, & \delta \end{pmatrix}$ mit der Determinante $+1$, durch die $(a, b, c)$ in $(a', b', c')$ übergeht.

Es ist dann
\weq{7}{
\begin{aligned}
a' &= a\alpha^2 + b\alpha\gamma + c\gamma^2,\\
b' &= 2a\alpha\beta + b(\alpha\delta + \beta\gamma) + c\gamma\delta,\\
c' &= a\beta^2 + b\beta\delta + c\delta^2,
\end{aligned}
}
\weq{8}{
\omega = \frac{\alpha\omega' + \beta}{\gamma\omega' + \delta}, \quad \omega' = \frac{\delta\omega - \beta}{-\gamma\omega + \alpha}\,,
}
und eine kleine Rechnung zeigt, daß $\omega$ und $\omega'$ entsprechende Wurzeln, d.\,h. Wurzeln mit dem gleichen Vorzeichen von $\sqrt{\Delta}$ sind.

Demnach erhalten wir für die zu $\varphi$ und $\varphi'$ gehörigen Ideale $\mathfrak{a}, \mathfrak{a}'$ zwei Basisformen:
\[
\begin{aligned}
\lambda &= a(x - \omega y),\\
\lambda' &= a'(x' - \omega'y'),
\end{aligned}
\]
und nach (8) ist
\[
\lambda = \frac{a[(\delta x - \beta y) + (\gamma x - \alpha y)\omega']}{\gamma\omega' + \delta}\,.
\]

Setzt man also
\[
x' = \delta x - \beta y, \quad y' = -\gamma x + \alpha y,
\]
so folgt
\[
\lambda = \frac{a}{a'(\gamma\omega' + \delta)}\,\lambda',
\]
da $a : a'(\gamma\omega' + \delta)$ eine Zahl in $\Omega$ ist, so folgt, daß $\mathfrak{a}$ und $\mathfrak{a}'$ äquivalent sind.

Wir haben also wieder in Übereinstimmung mit dem allgemeinen Satz \S\ 88, 4. und zugleich in näherer Bestimmung:

6. \gesperrt{Sind $(a, b, c)$ $(a', b', c')$ zwei äquivalente Formen mit positiven ersten Koeffizienten, und sind $\omega, \omega'$ entsprechende Wurzeln dieser Formen, so sind die Ideale $(a, a\omega)$ und $(a', a'\omega')$ äquivalent.}

Bei negativer Diskriminante haben die ersten Koeffizienten äquivalenter Formen immer dasselbe Vorzeichen, und man berücksichtigt nur die Formen, in denen dieses Vorzeichen positiv ist.
%CONT>>
